\scanpage{37}

Alors

\[
f: A\unsure{\longrightarrow|\mathbin{-}} E,\qquad u\longmapsto u.e
\]

est un iso respectant les structures quadratiques

\[
Q(f(u))=Q(u.e)=N(u)\underbrace{Q(e)}_{=1}=N(u)
\]

% slanted
\margin{--- Plus généralement, $Q(u.x)=N(u)Q(x)$ $(u\in\Gamma A,\ x\in\Gamma E)$}

par quoi la structure d'alg. de $A$ est transformée en l'unique str. d'Alg \inserted{quadratique} sur $\underline O$, admettant $e$ pour section unité, et $Q$ comme fonction quadratique Norme. La symétrie cherchée $\sigma_D$ peut être définie, par transport de structure, comme déduite de l'inv. canonique $\sigma$ de $A$

\[
\sigma:A\longrightarrow A\qquad \sigma(u)=\operatorname{Tr}.u-u.
\]

Donc on aura, par définition (posant $\sigma_e=\sigma_{\underline Oe}=\sigma_D$)

\[
\sigma_{\deleted{\unsure{D}}e}(ue)=\bar u.e
\]

On va donner une description équivalente, par

\[
\sigma_e(x)=\varphi_Q(e,x)e-x
\]

où $\varphi_Q$ est la forme bil. \inserted{sym} associée à $Q$. Cette formule ne fait que traduire celle pour \deleted{$A$}, $\sigma(u)=\operatorname{Tr}(u)-u$, compte tenu que \unsure{$\{$}

\[
\operatorname{Tr}(u)=\varphi_N(u,1).
\]

Avec la notation du ``produit scalaire'' (quand $2$ est inv sur $X$), on aura $\varphi_Q(e,x)=2e.x$, et on peut écrire

\[
\sigma_e(\deleted{\unsure{u}}x\deleted{\unsure{.e}})=2(e.x)e-x
\]

% archive sheet number: 36
